\documentclass[10pt,a4paper]{amsart}
\usepackage{amsmath,amssymb,amsthm,graphicx}
% Retained minimal preamble for the Stacks Project chapter more-algebra.tex (More on Algebra).
% The chapter file pinned-original-source.tex carries only \input{preamble} on its first line, so no
% preamble text lives inside the chapter itself.  The Stacks Project uses one single shared preamble for
% every chapter; the snapshot stages it identically as sources/stacks-chow/preamble.tex and
% sources/stacks-etale-cohomology/preamble.tex.  The block below is copied mechanically, line for line,
% from sources/stacks-chow/preamble.tex lines 14-18 (the Xy-pic setup) and 154-217 (the native theorem
% declarations together with the mathematical macros).  Nothing was added, removed or redefined; the
% wrapper class and the amsmath/amssymb/amsthm/graphicx base packages are supplied by the assembler.
% Chapter: more-algebra.tex -- More on Algebra.
% For commutative diagrams we use Xy-pic
\usepackage[all]{xy}

% We use 2cell for 2-commutative diagrams.
\xyoption{2cell}

% Theorem environments.
%
\theoremstyle{plain}
\newtheorem{theorem}[subsection]{Theorem}
\newtheorem{proposition}[subsection]{Proposition}
\newtheorem{lemma}[subsection]{Lemma}

\theoremstyle{definition}
\newtheorem{definition}[subsection]{Definition}
\newtheorem{example}[subsection]{Example}
\newtheorem{exercise}[subsection]{Exercise}
\newtheorem{situation}[subsection]{Situation}

\theoremstyle{remark}
\newtheorem{remark}[subsection]{Remark}
\newtheorem{remarks}[subsection]{Remarks}

\numberwithin{equation}{subsection}

% Macros
%
\def\lim{\mathop{\mathrm{lim}}\nolimits}
\def\colim{\mathop{\mathrm{colim}}\nolimits}
\def\Spec{\mathop{\mathrm{Spec}}}
\def\Hom{\mathop{\mathrm{Hom}}\nolimits}
\def\Ext{\mathop{\mathrm{Ext}}\nolimits}
\def\SheafHom{\mathop{\mathcal{H}\!\mathit{om}}\nolimits}
\def\SheafExt{\mathop{\mathcal{E}\!\mathit{xt}}\nolimits}
\def\Sch{\mathit{Sch}}
\def\Mor{\mathop{\mathrm{Mor}}\nolimits}
\def\Ob{\mathop{\mathrm{Ob}}\nolimits}
\def\Sh{\mathop{\mathit{Sh}}\nolimits}
\def\NL{\mathop{N\!L}\nolimits}
\def\CH{\mathop{\mathrm{CH}}\nolimits}
\def\proetale{{pro\text{-}\acute{e}tale}}
\def\etale{{\acute{e}tale}}
\def\QCoh{\mathit{QCoh}}
\def\Ker{\mathop{\mathrm{Ker}}}
\def\Im{\mathop{\mathrm{Im}}}
\def\Coker{\mathop{\mathrm{Coker}}}
\def\Coim{\mathop{\mathrm{Coim}}}

% Boxtimes
%
\DeclareMathSymbol{\boxtimes}{\mathbin}{AMSa}{"02}

%
% Macros for moduli stacks/spaces
%
\def\QCohstack{\mathcal{QC}\!\mathit{oh}}
\def\Cohstack{\mathcal{C}\!\mathit{oh}}
\def\Spacesstack{\mathcal{S}\!\mathit{paces}}
\def\Quotfunctor{\mathrm{Quot}}
\def\Hilbfunctor{\mathrm{Hilb}}
\def\Curvesstack{\mathcal{C}\!\mathit{urves}}
\def\Polarizedstack{\mathcal{P}\!\mathit{olarized}}
\def\Complexesstack{\mathcal{C}\!\mathit{omplexes}}
% \Pic is the operator that assigns to X its picard group, usage \Pic(X)
% \Picardstack_{X/B} denotes the Picard stack of X over B
% \Picardfunctor_{X/B} denotes the Picard functor of X over B
\def\Pic{\mathop{\mathrm{Pic}}\nolimits}
\def\Picardstack{\mathcal{P}\!\mathit{ic}}
\def\Picardfunctor{\mathrm{Pic}}
\def\Deformationcategory{\mathcal{D}\!\mathit{ef}}
\usepackage{cleveref}
\crefname{originallemma}{original lemma}{original lemmas}
\Crefname{originallemma}{Original lemma}{Original lemmas}
\makeatletter
\newlabel{more-algebra-lemma-CRT-henselization}{{0H7Q}{source}{}{sourcepointer-1}{}}
\newlabel{more-algebra-lemma-CRT-henselization@cref}{{[originallemma][0][]0H7Q}{[1][0][]source}{}{}{}}
\makeatother
\crefname{originallemma}{original lemma}{original lemmas}
\Crefname{originallemma}{Original lemma}{Original lemmas}
\makeatletter
\newlabel{lemma-henselization}{{0A02}{source}{}{sourcepointer-2}{}}
\newlabel{lemma-henselization@cref}{{[originallemma][0][]0A02}{[1][0][]source}{}{}{}}
\makeatother
\crefname{originallemma}{original lemma}{original lemmas}
\Crefname{originallemma}{Original lemma}{Original lemmas}
\makeatletter
\newlabel{categories-lemma-cofinal}{{04E7}{source}{}{sourcepointer-3}{}}
\newlabel{categories-lemma-cofinal@cref}{{[originallemma][0][]04E7}{[1][0][]source}{}{}{}}
\makeatother
\crefname{originallemma}{original lemma}{original lemmas}
\Crefname{originallemma}{Original lemma}{Original lemmas}
\makeatletter
\newlabel{algebra-lemma-flat-intersect-ideals}{{0BBY}{source}{}{sourcepointer-4}{}}
\newlabel{algebra-lemma-flat-intersect-ideals@cref}{{[originallemma][0][]0BBY}{[1][0][]source}{}{}{}}
\makeatother
\crefname{originallemma}{original lemma}{original lemmas}
\Crefname{originallemma}{Original lemma}{Original lemmas}
\makeatletter
\newlabel{algebra-lemma-chinese-remainder}{{00DT}{source}{}{sourcepointer-5}{}}
\newlabel{algebra-lemma-chinese-remainder@cref}{{[originallemma][0][]00DT}{[1][0][]source}{}{}{}}
\makeatother
\crefname{originallemma}{original lemma}{original lemmas}
\Crefname{originallemma}{Original lemma}{Original lemmas}
\makeatletter
\newlabel{categories-definition-cofinal}{{04E6}{source}{}{sourcepointer-6}{}}
\newlabel{categories-definition-cofinal@cref}{{[originallemma][0][]04E6}{[1][0][]source}{}{}{}}
\makeatother
\title{The Chinese remainder decomposition and henselization}
\author{The Stacks Project Authors}
\date{stacks-project:more-algebra.tex:local-snapshot-20261012; GNU FDL 1.2 or later}
\begin{document}
\maketitle
\setcounter{subsection}{1}
\noindent\textit{Editorial note.} All delivered mathematical text is the author's own, copied from the pinned Stacks Project chapter file more-algebra.tex; no statement or proof step was rewritten, reordered or removed.\par\medskip
\noindent\textit{Editorial note.} The chapter file pinned-original-source.tex carries only \textbackslash{}input\{preamble\} on its first line, so no preamble text lives inside the chapter. The retained file sources/stacks-more-algebra/retained-preamble.tex is a mechanical, line-for-line copy of the project's single shared preamble as staged in the snapshot (Xy-pic setup plus theorem declarations and mathematical macros). Nothing was added, removed or redefined, and the excerpt serves only as preamble, never as a span coordinate.\par\medskip
\noindent\textit{Editorial note.} Both counted spans start at the \textbackslash{}begin line of their environment, so the statement and the author's proof are delivered whole; surrounding prose of the section is deliberately outside the counted spans and is carried as standard prerequisites instead.\par\medskip
\noindent\textit{Editorial note.} The citation fields (source\_id, source\_version, authors, source\_url, license.name, source\_title) are identical to those of every other card this lane has produced from this chapter, so the batch selector treats them as one source.\par\medskip
\noindent\textit{Allowed standard prerequisites (editorial).}\begin{itemize}
\item Henselization of a pair (A, I) and its universal property. defined and discussed in the section, source lines 3180-3250, immediately before the counted statement. Supplies the object whose behaviour under a product decomposition is asserted.
\item Chinese remainder theorem for a ring and two ideals generating the unit ideal. used at several points of the proof, notably around proof lines 3290-3310. Provides the product decomposition of the rings involved.
\item The characterisation of a henselian pair through the lifting of factorisations of monic polynomials. used in the proof to verify the henselian property of the decomposed rings. Gives the property that has to be checked for each factor.
\end{itemize}\medskip
\ifdefined\hypertarget\hypertarget{sourcepointer-1}{}\fi
\noindent\textit{Source reference.} Original lemma 0H7Q: focused-seeds/AGP/sources/stacks-snapshot-20261011/tags.txt, source lines 21465--21465. The lemma's own canonical tag, quoted so that readers can resolve the label inside the delivered text; it points at the counted statement itself and carries no proof content.\par\medskip
\ifdefined\hypertarget\hypertarget{sourcepointer-2}{}\fi
\noindent\textit{Source reference.} Original lemma 0A02: focused-seeds/AGP/sources/stacks-snapshot-20261011/tags.txt, source lines 12620--12620. Navigation pointer to a statement the counted proof refers to but does not contain; nothing load-bearing is delegated to it.\par\medskip
\ifdefined\hypertarget\hypertarget{sourcepointer-3}{}\fi
\noindent\textit{Source reference.} Original lemma 04E7: focused-seeds/AGP/sources/stacks-snapshot-20261011/tags.txt, source lines 5753--5753. Navigation pointer to a statement the counted proof refers to but does not contain; nothing load-bearing is delegated to it.\par\medskip
\ifdefined\hypertarget\hypertarget{sourcepointer-4}{}\fi
\noindent\textit{Source reference.} Original lemma 0BBY: focused-seeds/AGP/sources/stacks-snapshot-20261011/tags.txt, source lines 14261--14261. Navigation pointer to a statement the counted proof refers to but does not contain; nothing load-bearing is delegated to it.\par\medskip
\ifdefined\hypertarget\hypertarget{sourcepointer-5}{}\fi
\noindent\textit{Source reference.} Original lemma 00DT: focused-seeds/AGP/sources/stacks-snapshot-20261011/tags.txt, source lines 509--509. Navigation pointer to a statement the counted proof refers to but does not contain; nothing load-bearing is delegated to it.\par\medskip
\ifdefined\hypertarget\hypertarget{sourcepointer-6}{}\fi
\noindent\textit{Source reference.} Original lemma 04E6: focused-seeds/AGP/sources/stacks-snapshot-20261011/tags.txt, source lines 5752--5752. Navigation pointer to a statement the counted proof refers to but does not contain; nothing load-bearing is delegated to it.\par\medskip
\begin{lemma}
\label{lemma-CRT-henselization}
Let $I_1, I_2, \dots, I_n \subset A$ be ideals in a ring $A$.
Suppose $I_i + I_j = A$ for $i \not = j$, $1 \leq i, j \leq n$.
Denote $A^h$ the henselization of the pair $(A, I_1 \cap \ldots \cap I_n)$
and $A_i^h$ the henselization of the pair $(A, I_i)$.
Then the natural morphism
$$
A^h \to \prod\nolimits_{i = 1, \ldots, n} A^h_i
$$
is an isomorphism.
\end{lemma}

\begin{proof}
The map comes from the functoriality of the construction.
If $n = 1$ the result is immediate. If $n > 1$, observe that
$(I_1 \cap \ldots \cap I_{n - 1}) + I_n = A$ by our condition
on the ideals. Thus if we prove the lemma for $n = 2$, then
the lemma follows for all $n$ by induction.

\medskip\noindent
The case $n = 2$. Choose $f_i \in I_i$ such that $f_1 + f_2 = 1$.
Recall that $A^h$ is constructed as the colimit of $A$-algebras $B$
such that $A \to B$ is \'etale and $A/I_1 \cap I_2 \to B/(I_1 \cap I_2)B$
is an isomorphism. See proof of Lemma \ref{lemma-henselization}.
Denote $Hens$ the category of such $A$-algebras $B$.
Similarly, denote $Hens_i$, $i = 1, 2$ the category of \'etale $A$-algebras
$B$ such that $A/I_i \to B/I_iB$ is an isomorphism.

\medskip\noindent
Consider the subcategory $Hens_1(f_2) \subset Hens_1$ consisting of
those $B$ such that $f_2$ is invertible. This is a cofinal subcategory
since if $B$ is in $Hens_1$, then $B_{f_2}$ is in $Hens_1$ as well
(hint: $B_{f_2}/I_1B_{f_2} = (B/I_1B)_{f_2} = B/I_1B$ because
$f_2$ is a unit modulo $I_1$). Similarly, the subcategory
$Hens_2(f_1) \subset Hens_2$ is cofinal. Hence to compute $A_1^h$ and
$A_2^h$ we may work with this smaller categories, see
Categories, Lemma \ref{categories-lemma-cofinal}.

\medskip\noindent
Consider the functor
$$
Hens_1(f_2) \times Hens_2(f_1) \to Hens,\quad (B_1, B_2) \mapsto B_1 \times B_2
$$
This functor makes sense as
\begin{align*}
(B_1 \times B_2)/(I_1 \cap I_2)(B_1 \times B_2)
& =
B_1/(I_1 \cap I_2)B_1 \times B_2/(I_1 \cap I_2)B_2 \\
& =
B_1/I_1B_1 \times B_2/I_2B_2 \\
& = A/I_1 \times A/I_2 \\
& = A/(I_1 \cap I_2)
\end{align*}
The second equality holds because
$(I_1 \cap I_2)B_1 = I_1B_1 \cap I_2B_1 = I_1B_1$ as $f_2 \in I_2$
is invertible in $B_1$
(we have also used Algebra, Lemma \ref{algebra-lemma-flat-intersect-ideals}).
The last equality holds by Algebra, Lemma \ref{algebra-lemma-chinese-remainder}.
To finish the proof, it suffices to prove that the displayed functor
is cofinal (see lemma quoted above).
The first condition of
Categories, Definition \ref{categories-definition-cofinal}
holds because an object $B$ of $Hens$ maps to $B_{f_2} \times B_{f_1}$
which is in the image of the functor. The second one always holds because
if we have $B \to B_1 \times B_2$ and $B \to B'_1 \times B'_2$ with
$B_1, B'_1 \in Hens_1(f_2)$ and $B_2, B'_2 \in Hens_2(f_1)$, then
we can set $B''_1 = B_1 \otimes_B B'_1$ which is in $Hens_1(f_2)$
and similarly $B''_2 = B_2 \otimes_B B''_2$ in $Hens_2(f_1)$
to find a map $B \to B''_1 \times B''_2$ fitting into the commutative diagram
$$
\xymatrix{
& B \ar[ld] \ar[d] \ar[rd] \\
B_1 \times B_2 \ar[r] & B_1'' \times B_2'' & B_1' \times B_2' \ar[l]
}
$$
This finishes the proof.
\end{proof}

\end{document}
